\documentclass[3p,times,procedia]{elsarticle}
\flushbottom
\usepackage{ecrc}
\usepackage{amsmath,amssymb}
\usepackage{booktabs}
\usepackage{graphicx}
\usepackage{float}
\usepackage{microtype}
\usepackage{etoolbox}
\usepackage[bookmarks=false]{hyperref}
\hypersetup{
  colorlinks,
  linkcolor=blue,
  citecolor=blue,
  urlcolor=blue,
  pdftitle={Continuous Behavioral Authentication for Mobile Devices Using an Adaptive LSTM-Based Intelligent Agent},
  pdfauthor={Vassiliy Serbin, Madi Gali, Maksim Iavich},
  pdfkeywords={continuous authentication; behavioral biometrics; LSTM; adaptive threshold; inertial sensors; HMOG; concept drift}
}
\volume{00}
\firstpage{1}
\journalname{Procedia Computer Science}
\runauth{Vassiliy Serbin et al.}
\jid{procs}
\makeatletter
\patchcmd{\ps@pprintTitle}{(2018)}{(2026)}{}{}
\patchcmd{\ps@headings}{(2018)}{(2026)}{}{}
\patchcmd{\ps@headings}{(2018)}{(2026)}{}{}
\pagestyle{headings}
\gdef\@copyrightyear{2026}
\makeatother
\begin{document}
\begin{frontmatter}
\dochead{The 17th International Conference on Emerging Ubiquitous Systems and Pervasive Networks\\
(EUSPN 2026), October 28--30, 2026, Almaty, Kazakhstan}

\title{Continuous Behavioral Authentication for Mobile Devices
Using an Adaptive LSTM-Based Intelligent Agent}

\author[a]{Vassiliy Serbin}
\author[b]{Madi Gali\corref{cor1}}
\ead{Gali.M@stud.satbayev.university}
\author[c]{Maksim Iavich}
\address[a]{Department of Information Systems, Institute of Automation and
Information Technologies, Satbayev University, Almaty 050013, Kazakhstan}
\address[b]{Department of Cybersecurity, Institute of Automation and
Information Technologies, Satbayev University, Almaty 050013, Kazakhstan}
\address[c]{Caucasus School of Technologies, Caucasus University, Tbilisi, Georgia}
\cortext[cor1]{Corresponding author.}

\begin{abstract}
Continuous authentication from smartphone inertial sensors must remain useful
when a legitimate user's score distribution changes, yet an adaptive decision
rule can also increase false acceptance or become vulnerable to profile
poisoning. We evaluate a two-layer long short-term memory (LSTM) classifier and
a lightweight agent that updates exponentially weighted first and second
moments of accepted confidence scores. The study uses accelerometer and
gyroscope data from 100 HMOG subjects. To prevent leakage, the protocol reads
only the physical sensor columns, constructs non-overlapping sequences within individual
sessions, keeps complete sessions disjoint between model fitting, validation,
and test, and estimates normalization and initial operating points without test
data. Linear SVM, random forest, and LSTM receive the same 10-second input and
achieve undrifted-session macro equal error rates of $29.92\%\pm0.80\%$,
$24.05\%\pm0.86\%$, and $21.87\%\pm1.01\%$, respectively. At the validation-calibrated
operating point, the frozen LSTM threshold and the EMA agent obtain HTERs of
$23.19\%\pm0.95\%$ and $24.53\%\pm1.28\%$ in a balanced alternating
zero-effort impostor stress test. At a final standardized-feature drift magnitude
$\lVert\delta\rVert_2=2$, their HTERs are $24.89\%\pm1.37\%$ and
$26.39\%\pm2.12\%$. Thus the tested EMA rule does not improve the frozen
threshold. These results provide a controlled ablation of the tested adaptive
component. The contribution is a reproducible, session-disjoint evaluation of
one prespecified EMA update rule, not a claim that adaptation is generally
beneficial.
\end{abstract}

\begin{keyword}
continuous authentication \sep behavioral biometrics \sep LSTM \sep
adaptive threshold \sep inertial sensors \sep HMOG \sep concept drift
\end{keyword}
\end{frontmatter}

\section{Introduction}
\label{sec:intro}

One-time authentication cannot protect an unlocked smartphone after the device
passes into another person's control. Continuous authentication (CA) instead evaluates behavioral
signals throughout a session, often without explicit interaction
\cite{alzubaidi_authentication_2016,stylios_behavioral_2021}. Accelerometers and
gyroscopes are attractive for this purpose because they are widely available
and capture user-specific handling and motion patterns
\cite{sitova_hmog_2016,abuhamad_sensorbased_2021}.

The difficulty is temporal change. A classifier trained during enrollment can
become poorly calibrated as behavior and activity context change
\cite{verma_data-driven_2022}. Retraining a deep model after every observation
is impractical on a mobile device, whereas a score-level update is inexpensive.
That convenience is not evidence of effectiveness: an update rule must reduce
false rejection without exposing the operating point to impostors, and it must
be evaluated on held-out data with all parameters fixed before test scoring.

We isolate the adaptation claim with a session-disjoint three-way protocol, an
explicit EMA decision equation whose parameters are fixed before testing, and
a paired frozen-LSTM versus LSTM+agent ablation at the same initial operating
point. Genuine attempts are processed in session and within-session order in
label-blind mixed streams. We report FAR, FRR, and HTER on undrifted sessions and
under synthetic drift, while limiting conclusions to the tested protocol.

\section{Related Work and Study Positioning}
\label{sec:related}

CA has a substantial literature. Broad surveys cover system components,
modalities, metrics, privacy, and open deployment problems
\cite{stylios_behavioral_2021,gonzalezmanzano_leveraging_2019,
ayeswarya_survey_2019}. A contemporary smartphone-focused survey catalogs
more than 140 sensor and behavioral-biometric approaches
\cite{abuhamad_sensorbased_2021}. The present work therefore does not claim that
continuous authentication, inertial sensing, LSTM classification, or EMA by
itself is new.

HMOG introduced hand-movement, orientation, and grasp features for smartphone
authentication \cite{sitova_hmog_2016}. Subsequent work applied augmentation,
deep sensor models, and adaptive verification to mobile behavioral data
\cite{verma_data-driven_2022,li_using_2019,abuhamad_autosen_2020}. LSTM and
hybrid recurrent networks are well established for sequential inertial
biometrics \cite{hochreiter_long_1997,zeng_gait-based_2021,giorgi_using_2021}.
Recent work has also studied smartphone authentication from gyroscope and
accelerometer data for mobile educational resources
\cite{serbin_behavioral_2026}. That study uses correlation-based decisions
rather than recurrent classification or online threshold adaptation and
therefore does not establish the performance of the present method.
A recent systematic review synthesizes 80 behavioral-biometric CA studies and
maps security, privacy, intelligent-agent, and deployment gaps
\cite{gali_behavioral_2026}.
Beyond CA, members of the author team have applied machine learning to
financial-feature analysis \cite{rakhmetulayeva_identifying_2025,
bolshibayeva_diskbased_2025,kulbayeva_data_processing_2024}, medical imaging
\cite{tyunina_cataracts_2026}, and UAV recognition \cite{semenyuk_uav_2024}.
These domain-specific studies
provide broader feature-engineering and deployment context; they are not used
as evidence for authentication performance.
Exponentially weighted statistics are likewise standard tools for tracking
non-stationary streams \cite{ross_exponentially_2012,gama_survey_2014}.

\begin{table}[H]
\centering
\small
\caption{Positioning of the present evaluation. ``Online update'' means a
label-blind state change during the chronological test stream; ``drift test''
means an explicit time-varying shift evaluated against a frozen operating point.
A check mark indicates that the cited work meets the stated definition.
Protocols and datasets differ, so error rates are not directly comparable.}
\label{tab:position}
\begin{tabular}{lcccc}
\toprule
Work & Inertial & Sequence model & Online update & Drift test \\
\midrule
Sitov\'{a} et al. \cite{sitova_hmog_2016} & $\checkmark$ & -- & -- & -- \\
AUToSen \cite{abuhamad_autosen_2020} & $\checkmark$ & $\checkmark$ & -- & -- \\
Verma et al. \cite{verma_data-driven_2022} & $\checkmark$ & -- & -- & -- \\
This study & $\checkmark$ & $\checkmark$ & $\checkmark$ & $\checkmark$ \\
\bottomrule
\end{tabular}
\end{table}

The methodological contribution is a paired evaluation: frozen and EMA-updated
LSTM thresholds start from the same validation profile and process identical
ordered test streams within each drift condition. This isolates the final-stage
update from classifier performance.

\section{Dataset and Leakage-Resistant Protocol}
\label{sec:data}

\subsection{HMOG data and parsing}

HMOG nominally contains 24 smartphone-interaction sessions for each of 100 participants,
including sitting and walking tasks \cite{sitova_hmog_2016}. We use the
accelerometer and gyroscope. An HMOG inertial CSV row contains metadata before
the three physical axes; consequently, only columns 4--6 are read as $x,y,z$.
This detail is leakage-critical because a session identifier must never enter
the feature vector. Accelerometer and gyroscope samples are paired by acquisition
order after truncation to their common length; no timestamp resampling is used.
Six sessions containing an empty inertial file for one subject are skipped, and
session 24 is absent for two other subjects; all retain enough valid sessions to
appear in every partition.

At the nominal 100-Hz sampling rate, the first five seconds are discarded. To
equalize session contributions and bound computation, the next 20 non-overlapping
windows of 100 samples are retained. Each window is represented by mean,
standard deviation, minimum, maximum, root mean square, and skewness for each of
six channels, producing 36 features. Ten adjacent windows form a 10-second
sequence. The sequence step is also ten windows, so evaluation endpoints do not
share raw samples. Sequences never cross a session boundary. The procedure yields
4,784 sequences from 2,392 valid sessions.

\subsection{Three-way session split}

For session number $r$, define $g(r)=(r-1)\bmod 5$. In outer fold $f$, sessions
with $g(r)=f$ are test sessions, sessions with
$g(r)=(f+1)\bmod 5$ are validation/calibration sessions, and the remaining
sessions are used to fit the model. Whole sessions, rather than sequences, are
assigned to partitions. Every included subject is required to occur in all
three partitions. This is a session-disjoint cross-validation protocol, not a
claim of leave-one-subject-out or prospective unseen-subject recognition.

A standard scaler is fitted only to the model-fitting sequences and then
applied unchanged to validation and test sequences. Early stopping and initial
threshold estimation use validation data; test labels are not used for either.
The same validation fold is used for checkpoint selection and claimant-specific
threshold/profile initialization; it is test-disjoint but not an independent
calibration set. This is closed-set verification over 100 enrolled identities.
For claimant $u$, the class-$u$ softmax probability is the score; sequences from
$u$ are genuine and those from the other 99 identities are zero-effort
impostors. All identities occur during fitting; unseen impostors are not tested.

\section{Method}
\label{sec:method}

\subsection{System overview}

The system has five stages:

\begin{center}
\fbox{sensor samples} $\rightarrow$
\fbox{1-s windows} $\rightarrow$
\fbox{36 features} $\rightarrow$
\fbox{LSTM score} $\rightarrow$
\fbox{frozen or EMA-updated threshold}
\end{center}

The first four stages are identical in the ablation. Only the last-stage
threshold update changes. Within each paired condition, FAR/FRR differences
therefore reflect the update rule rather than a different classifier or window.

\subsection{LSTM classifier and baselines}

The input is $\mathbf{X}\in\mathbb{R}^{10\times36}$. A two-layer LSTM with 64
hidden units per layer returns its last hidden state. Layer normalization,
a 32-unit ReLU layer, dropout of 0.3, and a 100-class softmax produce claimant
scores $s_{u,n}\in[0,1]$. Training uses cross-entropy and Adam
with learning rate $10^{-3}$, batch size 256, a maximum of 20 epochs, and
early-stopping patience 4.
Here ``agent'' denotes only the stateful EMA decision rule described below; it
does not retrain the LSTM or perform planning.

Linear SVM and random forest baselines receive the same sequence flattened to
360 dimensions. Linear SVM uses $C=0.1$; the forest contains 300
unrestricted-depth trees, uses square-root feature subsampling, and has minimum
leaf size 2. These fixed reference baselines compare classifier families at the
same temporal endpoint; they are neither tuned as state-of-the-art comparators
nor presented as adaptive systems.

\subsection{EMA adaptive agent}

For claimant $u$, the validation genuine scores initialize the first and second
moments $p_{u,0}$ and $q_{u,0}$. Before attempt $n$, the agent computes

\begin{align}
\sigma_{u,n-1} &= \sqrt{\max(q_{u,n-1}-p_{u,n-1}^{2},\epsilon^2)},\\
\tau_{u,n} &= \max(0,p_{u,n-1}-k_u\sigma_{u,n-1}),\\
d_{u,n} &= \mathbb{I}[s_{u,n}\geq\tau_{u,n}].
\end{align}

\begin{samepage}
When $d_{u,n}=1$, the accepted score updates both moments:

\begin{align}
p_{u,n} &= (1-\alpha)p_{u,n-1}+\alpha s_{u,n},\\
q_{u,n} &= (1-\alpha)q_{u,n-1}+\alpha s_{u,n}^{2}.
\end{align}
\end{samepage}

Here $\epsilon=10^{-4}$ supplies a strictly positive scale floor when
validation-score variance is negligible. After a rejection the state is unchanged. Thus repeated rejected events
cannot lower the threshold. An accepted impostor can still poison the profile;
the evaluation deliberately preserves that risk instead of using ground-truth
labels to gate updates.

We fix $\alpha=0.05$ a priori, corresponding to a characteristic update horizon
of approximately $1/\alpha=20$ accepted scores. For each claimant, the validation genuine and
impostor scores determine an EER threshold $\tau_{u,0}$; the corresponding
$k_u=(p_{u,0}-\tau_{u,0})/\sigma_{u,0}$ makes the first adaptive threshold
exactly equal to that test-independent, validation-calibrated operating point. Thus $k_u$ is
not selected on test scores. The frozen LSTM ablation uses the same
$\tau_{u,0}$ but never updates it. Decisions use the current claimant score; no
additional score buffer is used.

\section{Evaluation}
\label{sec:evaluation}

For each claimant, EER is approximated at the held-out ROC point minimizing
$|\mathrm{FAR}-\mathrm{FRR}|$. Claimant values are macro-averaged within each
fold, after which mean and standard deviation are computed over five folds.
This test-derived point summarizes separability and is never used for decisions.

The frozen-threshold and EMA-agent decisions start from the same
validation-calibrated operating point. Genuine test attempts follow session and within-session
chronology. Each is alternated with a zero-effort impostor score sampled without
replacement uniformly from all test sequences of the other 99 identities;
sampling is repeated ten times, with the validation profile
reinitialized before each repeat. Genuine chronology is unchanged while the
sampled impostor stream varies. This balanced 50\% impostor stream is
an explicit stress test, not an estimate of real attack prevalence. Because the
agent sees no labels, an accepted impostor updates its profile. We report macro
false acceptance rate (FAR), false rejection rate (FRR), and
$\mathrm{HTER}=(\mathrm{FAR}+\mathrm{FRR})/2$.

To study score drift, one unit direction sampled from $\mathcal{N}(0,I_{36})$
with seed $42+f$ is ramped linearly over each subject's test chronology to final
$\ell_2$ magnitudes $\lVert\delta\rVert_2\in\{0,1,2\}$. At each endpoint, the
same 36-dimensional shift is added to all ten time steps of a sequence; the
stated norm refers to the per-time-step standardized feature vector. The
transformation is applied to all test samples, including samples subsequently
drawn as impostors. The same trained model is used. This is a synthetic
feature-shift sensitivity analysis, not a model of natural behavioral drift;
only frozen-versus-EMA comparisons within each level are paired. Impostor
sampling uses seed $42+100f+10\lVert\delta\rVert_2$.

\section{Results}
\label{sec:results}

\begin{table}[H]
\centering
\caption{Macro EER on undrifted held-out sessions (mean $\pm$ standard deviation
over five outer folds). All methods use the same 10-second sequences.}
\label{tab:clean}
\begin{tabular}{lc}
\toprule
Model & EER (\%) \\
\midrule
Linear SVM & $29.92\pm0.80$ \\
Random forest & $24.05\pm0.86$ \\
LSTM (softmax score) & $21.87\pm1.01$ \\
\bottomrule
\end{tabular}
\end{table}

Table~\ref{tab:clean} compares score separability. Table~\ref{tab:ablation}
isolates the agent at a threshold estimated without test data. The agent is not
assigned an additional post-hoc EER; its chronological, label-blind decisions are reported.

\begin{table}[H]
\centering
\small
\caption{Frozen-threshold and EMA-agent operating points in the balanced
alternating stream (mean $\pm$ standard deviation over five folds).}
\label{tab:ablation}
\begin{tabular}{llccc}
\toprule
$\lVert\delta\rVert_2$ & Decision rule & FAR (\%) & FRR (\%) & HTER (\%) \\
\midrule
0 & Frozen & $22.55\pm1.13$ & $23.82\pm2.27$ & $23.19\pm0.95$ \\
0 & EMA agent & $24.46\pm2.12$ & $24.60\pm2.04$ & $24.53\pm1.28$ \\
1 & Frozen & $22.43\pm0.69$ & $24.82\pm1.97$ & $23.63\pm1.10$ \\
1 & EMA agent & $24.30\pm2.00$ & $25.73\pm2.64$ & $25.02\pm1.66$ \\
2 & Frozen & $22.40\pm0.85$ & $27.37\pm2.91$ & $24.89\pm1.37$ \\
2 & EMA agent & $24.53\pm2.11$ & $28.25\pm4.25$ & $26.39\pm2.12$ \\
\bottomrule
\end{tabular}
\end{table}

On undrifted data, the LSTM lowers macro EER by 2.18 percentage points relative to
the random forest and by 8.05 points relative to the linear SVM.

The agent ablation is negative. Adaptation increases HTER by 1.34, 1.39, and
1.51 points at drift magnitudes 0, 1, and 2, respectively, using unrounded fold
means. The frozen threshold has lower HTER in all five folds at every level.
Both FAR and FRR rise on average, so this is not a deliberate security--usability
trade-off. Observed HTER is higher at larger nominal drift, although independent
impostor resampling also contributes to cross-level variation.
These fold summaries are descriptive: training and validation partitions
overlap across outer folds, so the reported standard deviations are not
confidence intervals and no statistical-significance claim is made.

\section{Discussion}
\label{sec:discussion}

Table~\ref{tab:clean} measures discrimination, whereas
Table~\ref{tab:ablation} evaluates the test-independent, validation-calibrated operating
point. Low EER alone does not establish an adaptation benefit.

The label-blind accepted-only update cannot lower the threshold after a rejection,
cannot adapt to sub-threshold genuine scores, and remains vulnerable to poisoning
after a false acceptance. Future deployment studies should therefore evaluate
high-confidence update gates, step-up authentication, and poisoning defenses.

\subsection{Limitations}

Evaluation is limited to HMOG; six empty-file sessions and two absent directories
are excluded. Row-order sensor pairing may leave temporal misalignment, and the
20-second excerpt may emphasize early-session behavior. The modulo split is not
a rolling-origin forecast, produces unequal fold sizes, and is not explicitly
stratified by activity or context. The alternating 50\% impostor stream fixes
genuine before impostor and does not model attack prevalence or bursts. Independent
impostor draws at each drift level confound cross-level changes with Monte Carlo
variation; only frozen-versus-EMA comparisons within a level are paired. Drift
uses one direction per fold. Only one prespecified EMA rate, $\alpha=0.05$, is
evaluated; the negative result does not generalize to other update rates or
adaptation rules. Future work requires independent datasets, natural sensor
shifts, confidence intervals accounting for subject/session clustering, and
on-device latency and energy tests.

\paragraph{Data and code availability} Code, protocol metadata, and fold results
are archived in the public
\href{https://github.com/madigali/continuous-mobile-authentication-hmog/releases/tag/v1.0.0}{GitHub release v1.0.0}.
HMOG data are not redistributed.

\section{Conclusion}

Under this session-disjoint HMOG protocol, the LSTM has the lowest undrifted
macro EER ($21.87\%\pm1.01\%$), but EMA adaptation raises HTER from $23.19\%$
to $24.53\%$ without drift and from $24.89\%$ to $26.39\%$ at drift magnitude 2.
The tested rule does not improve the validation-calibrated frozen threshold;
independent datasets, natural drift, poisoning, and on-device tests remain necessary.

\bibliographystyle{elsarticle-num}
\bibliography{references}
\end{document}
